%=============================================================================
% WAVE 4, ITERATION 14: P1 + P8 DEAD PATENTS --- AUDIT & NEW FINDINGS
% Neutrino Mass Existential Threat, DESI DR2 Sign-Change Audit,
% S8 Tension 2026 Review, Preferred-Frame PPN Constraint, Analog Fiber Update
%
% Agent: P1+P8 Dead Patents (W4i14) | Model: Opus 4.6
% Date: 20 March 2026
%
% MANDATE:
%   1. Confirm P1 and P8 remain dead (1 paragraph)
%   2. Audit and cross-check claims from W4i13 agents
%   3. Find things NOBODY ELSE will find
%   4. TOP 5 findings, derivation chains, flag inconsistencies
%
% INHERITED FACTS (not re-derived):
%   - P1 DEAD (W3i7 FINAL, 4th consecutive confirmation)
%   - P8 DEAD (W3i7 FINAL, 4th consecutive confirmation)
%   - All P1/P8-relevant bounds unchanged since W4i12
%   - W4i13 Finding 1 (analog gravity), Finding 2 (NANOGrav monopole),
%     Finding 3 (MAQRO), Finding 4 (phase transitions), Finding 5 (GRIN)
%     all remain valid.
%=============================================================================

\documentclass[11pt]{article}
\usepackage[margin=2.5cm]{geometry}
\usepackage{amsmath,amssymb,mathtools}
\usepackage{booktabs}
\usepackage{enumitem}
\usepackage{float}
\usepackage{xcolor}
\usepackage[most]{tcolorbox}
\usepackage{hyperref}

\newcommand{\ee}[1]{\times 10^{#1}}
\newcommand{\lamdiff}{|\lambda{-}1|}
\newcommand{\psifield}{\psi}
\newcommand{\order}[1]{\mathcal{O}(#1)}
\newcommand{\LCDM}{$\Lambda$CDM}
\newcommand{\DFD}{\textsc{dfd}}
\newcommand{\Geff}{G_{\rm eff}}
\newcommand{\Hzero}{H_0}

\title{\textbf{Wave 4 Iteration 14: P1 + P8 Dead Patents --- Audit Update}\\[0.5em]
\large Neutrino Mass Existential Threat $\cdot$ DESI DR2 Sign-Change Audit\\
$S_8$ Tension 2026 Review $\cdot$ Preferred-Frame PPN Test\\
Optical Fiber Analog Update}
\author{DFD Research Programme --- P1+P8 Agent, W4i14 (Opus 4.6)}
\date{20 March 2026}

\begin{document}
\maketitle

%=============================================================================
\section*{Executive Summary: TOP 5 FINDINGS (Ranked by Impact)}
%=============================================================================

\begin{tcolorbox}[colback=red!5!white, colframe=red!60!black,
  title=\textbf{TOP 5 FINDINGS --- W4i14}]

\textbf{Finding 1 (EXISTENTIAL --- AUDIT OF CORPUS NEUTRINO MASS):}
\textbf{DFD's $\Sigma m_\nu = 61.4$~meV prediction is now EXCLUDED at
$>$95\% CL by DESI DR2 + Planck PR4 + CMB lensing.}  The tightest
current cosmological bound (arXiv:2507.12401, Feldman--Cousins
frequentist analysis) gives $\Sigma m_\nu < 53$~meV at 95\% CL, which
is \emph{below} the normal-ordering floor of 59~meV.  DFD predicts
$\Sigma m_\nu = 61.4$~meV (Tier~1, zero-parameter), which exceeds this
bound by 8.4~meV.  The corpus margin was listed as ``2.8~meV vs DESI
DR2 bound'' (from the earlier 72~meV DESI DR1 bound).  \textbf{The
margin is now $-8.4$~meV --- DFD is on the wrong side.}  This is the
single most dangerous quantitative result in the programme.  Two
mitigating factors: (i) the 53~meV bound assumes \LCDM; under $w_0 w_a$CDM
the bound relaxes to $\sim 120$~meV (arXiv:2406.14554), and DFD's
dark-energy sector differs from \LCDM; (ii) the frequentist vs.\
Bayesian bounds differ --- the Bayesian posterior gives $\Sigma m_\nu <
72$~meV (95\% CL) under \LCDM, with which DFD's 61.4~meV is consistent.
\textbf{Nevertheless, the corpus must urgently flag that the DFD
Boltzmann code is needed to self-consistently derive the neutrino mass
bound \emph{within the DFD cosmology}, not under the \LCDM{} assumption.}

\textbf{Finding 2 (CRITICAL --- AUDIT OF W4i13 SURVEY AGENT):}
\textbf{The W4i12 $D_H/r_d$ sign-change table has a SIGN ERROR at
$z = 2.33$ that was flagged by the W4i13 Survey Agent but remains
unresolved.}  The W4i12 table lists $D_H^{\rm DFD}/r_d = 8.62$ at
$z = 2.33$, which is \emph{above} \LCDM{} ($= 8.52$), a $+1.2\%$
deviation.  But the DFD sign-change prediction requires $D_H^{\rm DFD}
< D_H^{\LCDM}$ at $z > 1.78$.  The published DESI DR2 Lyman-$\alpha$
value $8.632 \pm 0.101$ is also above \LCDM.  There are two
possibilities: (a) the DFD sign-change prediction is qualitatively
correct but the W4i12 table has a computational error at $z = 2.33$
(the ``50$\times$ amplitude error'' noted in iteration context), or
(b) the sign change occurs at $z > 2.33$, not $z = 1.78$.  Either way,
\textbf{the W4i12 $D_H/r_d$ table cannot be used for observer
predictions until the Boltzmann code produces self-consistent values}.
The corpus should downgrade this from ``Tier 1 prediction'' to
``Tier 2 (pending numerical validation)''.

\textbf{Finding 3 (HIGH IMPACT --- NEW):}
\textbf{DFD's preferred foliation predicts a non-zero PPN
$\alpha_1$ parameter, testable by Lunar Laser Ranging.}
DFD has a preferred frame (absolute simultaneity, constant-$t$
surfaces; section02\_formalism.tex \S2.6).  In PPN theory, a preferred
frame generates non-zero $\alpha_1$ and $\alpha_2$ parameters.  For
DFD's scalar sector, the preferred-frame effect arises from the
$|\nabla\psi|$ dependence of $\mu(x)$: the gradient is frame-dependent
at $\order{v^2/c^2}$.  The leading PPN preferred-frame contribution is:
\begin{equation}
\alpha_1^{\rm DFD} = \frac{4 v_{\odot}^2}{c^2}\frac{\mu'(x_{\odot})}
{\mu(x_{\odot})}\bigg|_{x_{\odot} \approx 10^{7}},
\end{equation}
where $v_\odot \approx 370$~km/s is the solar system velocity relative
to the CMB.  In the Newtonian regime ($x_\odot \gg 1$), $\mu(x) =
x/(1+x)$ gives $\mu'/\mu = 1/(x(1+x)) \approx 1/x^2 \sim 10^{-14}$.
Thus:
\begin{equation}
\alpha_1^{\rm DFD} \sim 4 \times \frac{(3.7\ee{5})^2}{(3\ee{8})^2}
\times 10^{-14} = 4 \times 1.5\ee{-6} \times 10^{-14}
\sim 6\ee{-20}.
\end{equation}
\textbf{Derivation chain}: $v_\odot = 370$~km/s (CMB dipole) $\to$
$v^2/c^2 = 1.52\ee{-6}$ $\to$ $x_\odot = a_\odot/a_0 = GM_\odot/
(r_\oplus^2 a_0) = 5.3\ee{7}$ $\to$ $\mu'/\mu = 1/(x(1+x)) \approx
3.6\ee{-16}$ $\to$ $\alpha_1 = 4 \times 1.52\ee{-6} \times 3.6\ee{-16}
= 2.2\ee{-21}$.

Current LLR bound: $|\alpha_1| < 10^{-4}$ (Williams et al.\ 2004).
\textbf{DFD's prediction is 17 orders of magnitude below the current
bound.}  This confirms that DFD's preferred frame is undetectable by
any foreseeable experiment in the solar system, but the \emph{theoretical
prediction is non-zero} --- making DFD technically distinguishable from
Lorentz-invariant theories if precision ever reaches $10^{-20}$.  This
is a \textbf{new honest negative}: DFD violates Lorentz invariance at
an undetectably tiny level.

\textbf{Finding 4 (HIGH IMPACT --- AUDIT OF $S_8$ CLAIM):}
\textbf{The 2026 $S_8$ tension landscape has shifted: some probes now
agree with Planck, weakening DFD's $S_8$ prediction as a discriminator.}
A comprehensive 2026 review (arXiv:2602.12238) shows that HSC weak
lensing combined with spectroscopic data finds no evidence for $S_8$
tension with CMB.  Meanwhile, the DES Y6 $S_8 = 0.789$ (cited in
context as ``validating DFD'') uses angular scales where DFD predicts
$S_8 \sim 0.79$--$0.83$.  However, if the $S_8$ tension dissolves
entirely (i.e., all surveys converge on Planck's $S_8 \approx 0.83$),
DFD's scale-dependent $S_8$ prediction becomes merely consistent rather
than explanatory.  \textbf{The corpus should note that $S_8$ validation
is conditional on the tension persisting.}  Euclid's first cosmology
release (October 2026) will be decisive.

\textbf{Finding 5 (MODERATE --- NEW):}
\textbf{Nonlinear optical fiber analogs now have direct literature
support (arXiv:2512.15695, December 2025).}  A recent review on
nonlinear fiber optics and optical analogues to gravitational phenomena
(arXiv:2512.15695v1) systematically develops the correspondence between
graded-index fiber optics and curved spacetime, including the Gordon
optical metric formalism that DFD uses as its foundation.  This paper
establishes that the mathematical framework of the W4i13 ``psi-simulator''
concept (Finding 5) is well-grounded in current fiber-optics literature.
It also identifies specific fiber nonlinearities (Kerr effect, saturable
absorption) that can implement position-dependent refractive indices
matching the DFD $\mu$-crossover.  \textbf{The W4i13 GRIN fiber concept
now has direct literature backing and can be cited as an established
analog gravity technique.}

\end{tcolorbox}


%=============================================================================
\part{P1 and P8 Death Confirmation}
%=============================================================================

\section{P1 and P8: Confirmed Dead (4th Consecutive Iteration)}

P1 (drag reduction / Doppler sensing) and P8 (deep-space psi-communications)
remain irrecoverably dead.  The fundamental barriers are unchanged:
P1 requires $\lamdiff > 1.85\ee{-2}$ against the SQMS bound of
$1.56\ee{-9}$ (gap $> 10^7$), with the strongest laboratory signal at
$\sim 0.81$~zs ($10^{15}$ gap to detection).  P8 requires $1.6\,M_\oplus$
for 1~kbit/s at 1~AU.  No 2025--2026 developments in fifth-force
constraints, scalar-tensor gravity, or gravitational communication theory
change these assessments.  This is the 4th consecutive confirmation
since the W3i7 FINAL ruling.  \textbf{P1 and P8: DEAD, zero resources.}


%=============================================================================
\part{Detailed Findings}
%=============================================================================

%---------------------------------------------------------------------
\section{Finding 1: Neutrino Mass --- Existential Threat}
\label{sec:neutrino}
%---------------------------------------------------------------------

\subsection{The Problem}

The DFD corpus predicts $\Sigma m_\nu = 61.4$~meV as a Tier~1,
zero-parameter prediction (W4i10, verified W4i11).  The derivation
chain is:
\begin{equation}
\Sigma m_\nu^{\rm DFD} = \frac{3 \Hzero^2}{8\pi G}
\frac{\Omega_\nu^{\rm DFD}}{n_\nu}
\cdot \frac{1}{(T_\nu/T_\gamma)^3}
= 61.4 \;\text{meV},
\end{equation}
where $\Omega_\nu^{\rm DFD}$ is set by the DFD modified Friedmann
equation with $\Geff(k)$.

\subsection{The New Constraint}

The DESI DR2 + Planck PR4 + CMB lensing combination (arXiv:2507.12401,
July 2025) gives:
\begin{equation}
\Sigma m_\nu < 53 \;\text{meV} \quad (95\%\;\text{CL, Feldman--Cousins}).
\end{equation}
This is below the normal-ordering floor $\Sigma m_\nu^{\rm NO} \geq
59$~meV.  The DFD prediction of 61.4~meV is excluded at $> 95\%$ CL
under this analysis.

\subsection{Mitigating Factors}

\begin{enumerate}[nosep]
\item \textbf{Model dependence}: The 53~meV bound assumes \LCDM.  Under
  $w_0 w_a$CDM (which DESI DR2 itself prefers at $2.8$--$4.2\sigma$),
  the bound relaxes to $\sim 120$~meV.  DFD's cosmology is neither
  \LCDM{} nor $w_0 w_a$CDM, so the correct constraint must be derived
  within the DFD framework.

\item \textbf{Frequentist vs.\ Bayesian}: The Bayesian 95\% upper
  limit under \LCDM{} is $\Sigma m_\nu < 72$~meV, with which DFD's
  61.4~meV is consistent.  The frequentist Feldman--Cousins procedure
  gives the tighter 53~meV bound.

\item \textbf{Physical prior}: The normal-ordering floor (59~meV)
  provides a physical prior that is not included in the cosmological
  analysis.  Imposing this prior shifts the constraint upward.

\item \textbf{DFD Boltzmann code}: DFD's $\Geff(k)$ modifies the
  growth history, which changes the degeneracy between $\Sigma m_\nu$
  and other parameters.  The self-consistent DFD constraint on
  $\Sigma m_\nu$ could be significantly different from the \LCDM-derived
  value.
\end{enumerate}

\subsection{Assessment}

The neutrino mass constraint has transitioned from ``margin of 2.8~meV''
(corpus W4i12, citing the earlier 72~meV Bayesian \LCDM{} bound from
DESI DR1) to ``excluded by 8.4~meV'' under the tightest frequentist
analysis.  This is \textbf{not yet a clean falsification} because the
bound is model-dependent, but it is the \emph{closest any constraint has
come to falsifying a Tier~1 DFD prediction}.

\begin{tcolorbox}[colback=red!10, colframe=red!70!black]
\textbf{URGENT ACTION}: The DFD Boltzmann code (CLASS/CAMB modification)
must derive $\Sigma m_\nu$ constraints self-consistently within DFD
cosmology.  Until this is done, the neutrino mass prediction should be
flagged as \textbf{Tier 1 (AT RISK)} rather than validated.  The corpus
statement ``margin only 2.8~meV vs DESI DR2 bound'' is \textbf{STALE}
and must be updated to reflect the 53~meV frequentist bound.
\end{tcolorbox}


%---------------------------------------------------------------------
\section{Finding 2: $D_H/r_d$ Sign-Change Table Audit}
\label{sec:sign-change}
%---------------------------------------------------------------------

\subsection{The Inconsistency}

The W4i12 Boltzmann specification table lists $D_H/r_d$ predictions at
each DESI redshift bin.  The W4i13 Survey Agent (Finding 1 of
W4i13\_Survey\_update.tex) identified a critical internal inconsistency:

\begin{itemize}[nosep]
\item The DFD prediction claims $D_H^{\rm DFD} < D_H^{\LCDM}$ for
  $z > 1.78$.
\item The table lists $D_H^{\rm DFD}/r_d = 8.62$ at $z = 2.33$, which
  is $+1.2\%$ \emph{above} \LCDM{} ($= 8.52$).
\item These two statements are mutually contradictory.
\end{itemize}

\subsection{Root Cause}

The iteration context notes a ``50$\times$ amplitude error'' in the
W4i12 $D_H/r_d$ table.  The sign-change is described as ``qualitatively
robust but quantitatively wrong.''  This means the \emph{direction}
of the sign change (DFD predicts higher $D_H$ at low $z$, lower $D_H$
at high $z$) may be correct, but the \emph{crossover redshift} and
\emph{amplitude} are unreliable.

\subsection{DESI DR2 Data}

DESI DR2 Lyman-$\alpha$ at $z_{\rm eff} = 2.33$: $D_H/r_d = 8.632
\pm 0.101$.  This is \emph{above} \LCDM{} by $+1.3\%$, consistent
with no sign change at $z = 2.33$.  The low-$z$ DESI bins (LRG,
$z \sim 0.5$--$0.9$) show $D_H/r_d$ within $\pm 1\%$ of \LCDM.

\subsection{Assessment}

The sign-change prediction is currently \textbf{untestable with the
existing DFD computational tools} because the table values are known to
contain errors.  The DESI DR2 data are consistent with \LCDM{} at all
redshifts within $1$--$2\sigma$.  No sign change is visible in the
data, but the current precision ($\sim 1$--$3\%$ per bin) is
insufficient to rule out a $\sim 0.5$--$1\%$ DFD effect.

\begin{tcolorbox}[colback=orange!10, colframe=orange!70!black]
\textbf{RECOMMENDATION}: Downgrade the $D_H/r_d$ sign change from
``Tier 1 prediction'' to ``Tier 2 (pending Boltzmann code)''.  The
W4i12 table should be marked as \textbf{UNRELIABLE} in the corpus.
No observer paper should cite the specific numerical values until
the Boltzmann code produces self-consistent predictions.
\end{tcolorbox}


%---------------------------------------------------------------------
\section{Finding 3: Preferred-Frame PPN Parameter $\alpha_1$}
\label{sec:PPN}
%---------------------------------------------------------------------

\subsection{Theoretical Basis}

DFD explicitly posits a preferred foliation: absolute simultaneity on
constant-$t$ surfaces (section02\_formalism.tex, \S2.6).  In PPN
theory, this generates non-zero preferred-frame parameters $\alpha_1$,
$\alpha_2$, $\alpha_3$.

\subsection{Derivation}

In DFD, the scalar field equation involves $|\nabla\psi|/a_\star$, which
depends on the spatial gradient in the preferred frame.  A boost by
velocity $\mathbf{v}$ modifies the gradient:
\begin{equation}
|\nabla\psi|^2 \to |\nabla\psi|^2 - \frac{1}{c^2}
(\mathbf{v}\cdot\nabla\psi)^2 + \order{v^4/c^4}.
\end{equation}
The $\mu$-function's argument changes by:
\begin{equation}
\delta x = -\frac{v^2}{c^2}\frac{x\cos^2\theta}{1}
\sim -\frac{v^2}{c^2}x,
\end{equation}
where $\theta$ is the angle between $\mathbf{v}$ and $\nabla\psi$.
The PPN preferred-frame parameter is:
\begin{equation}
\alpha_1^{\rm DFD} = 4\frac{v_{\rm CMB}^2}{c^2}
\cdot \frac{d\ln\mu}{d\ln x}\bigg|_{x = x_\odot}.
\end{equation}

For $\mu(x) = x/(1+x)$:
\begin{equation}
\frac{d\ln\mu}{d\ln x} = \frac{x\mu'(x)}{\mu(x)}
= \frac{x \cdot 1/(1+x)^2}{x/(1+x)} = \frac{1}{1+x}.
\end{equation}

At the solar surface acceleration $x_\odot = a_\odot/a_0
= GM_\odot/(r_\oplus^2 a_0) = (6.67\ee{-11}\times 2\ee{30})/
((1.5\ee{11})^2 \times 1.12\ee{-10}) = 5.3\ee{7}$:
\begin{equation}
\frac{d\ln\mu}{d\ln x}\bigg|_{x_\odot} = \frac{1}{5.3\ee{7}}
= 1.9\ee{-8}.
\end{equation}

With $v_{\rm CMB} = 370$~km/s:
\begin{equation}
\alpha_1^{\rm DFD} = 4 \times \frac{(3.7\ee{5})^2}{(3\ee{8})^2}
\times 1.9\ee{-8}
= 4 \times 1.52\ee{-6} \times 1.9\ee{-8}
= 1.2\ee{-13}.
\end{equation}

\textbf{Derivation chain}:
\begin{itemize}[nosep]
\item $v_{\rm CMB} = 370$~km/s (Planck CMB dipole)
\item $v^2/c^2 = (3.7\ee{5}/3\ee{8})^2 = 1.52\ee{-6}$
\item $a_\odot = GM_\odot/r_\oplus^2 = 5.93\ee{-3}$~m/s$^2$
\item $x_\odot = a_\odot/a_0 = 5.93\ee{-3}/1.12\ee{-10} = 5.3\ee{7}$
\item $d\ln\mu/d\ln x = 1/(1+x) \approx 1/x = 1.9\ee{-8}$
\item $\alpha_1 = 4 \times 1.52\ee{-6} \times 1.9\ee{-8}
  = 1.2\ee{-13}$
\end{itemize}

Current best bound: $|\alpha_1| < 4\ee{-5}$ (LLR, Nordtvedt 2001).
DFD prediction: $\alpha_1 \sim 10^{-13}$, eight orders of magnitude
below the bound.

\subsection{Significance}

This is a \textbf{new honest negative}: DFD's preferred frame is
undetectable by any foreseeable experiment, confirming that the theory's
Lorentz violation is purely formal in the solar system.  However, it is
a genuine theoretical prediction that distinguishes DFD from
Lorentz-invariant scalar-tensor theories.  It should be added to the
corpus as a ``zero-parameter prediction (not testable within current
technology)''.


%---------------------------------------------------------------------
\section{Finding 4: $S_8$ Tension 2026 Landscape Audit}
\label{sec:S8}
%---------------------------------------------------------------------

\subsection{Current State}

A comprehensive 2026 review of the $S_8$ tension (arXiv:2602.12238)
reveals a shifting landscape:

\begin{center}
\begin{tabular}{llcc}
\toprule
\textbf{Survey} & \textbf{Method} & $S_8$ & \textbf{Tension with
Planck} \\
\midrule
Planck 2018 & CMB & $0.832 \pm 0.013$ & --- \\
DES Y6 & Cosmic shear & $0.789$ & $\sim 2$--$3\sigma$ \\
KiDS-Legacy & Cosmic shear & $0.815 \pm 0.016$ & $< 1\sigma$ \\
HSC + spectro & WL + spec-$z$ & $\sim 0.82$ & $< 0.5\sigma$ \\
eROSITA & Cluster counts & $0.86 \pm 0.01$ & $+2\sigma$ (high!) \\
\bottomrule
\end{tabular}
\end{center}

\subsection{Impact on DFD}

DFD predicts scale-dependent $S_8$: 0.73--0.76 at large scales
($\theta > 100'$), 0.81--0.83 at small scales ($\theta < 10'$).  This
is \emph{consistent} with a picture where large-scale surveys
(DES, KiDS) find lower $S_8$ and cluster counts (eROSITA) find higher
$S_8$.  However:

\begin{enumerate}[nosep]
\item If the tension dissolves and all probes converge on $S_8 \approx
  0.83$, DFD's large-scale prediction (0.73--0.76) would be in
  \textbf{tension at 3--4$\sigma$} with the data.
\item If the tension persists with a scale-dependent pattern, DFD is
  \textbf{strongly favored} as an explanation.
\end{enumerate}

\textbf{Assessment}: The DES Y6 $S_8 = 0.789$ ``validation'' of DFD
cited in the corpus is conditional.  The corpus should note: ``$S_8$
validation requires confirmation that the tension is scale-dependent,
not merely a systematic.  Euclid first cosmology release (October 2026)
is decisive.''


%---------------------------------------------------------------------
\section{Finding 5: Optical Fiber Analog --- Literature Update}
\label{sec:fiber}
%---------------------------------------------------------------------

\subsection{New Reference}

A December 2025 review paper (arXiv:2512.15695) systematically develops
nonlinear fiber optics as an analog gravity framework, using precisely
the Gordon optical metric formalism that underlies DFD.  Key points
relevant to the W4i13 GRIN fiber concept:

\begin{enumerate}[nosep]
\item The correspondence between graded-index fibers and curved
  spacetimes is established for both linear and nonlinear regimes.
\item Kerr-type nonlinearities ($n = n_0 + n_2 I$) provide
  intensity-dependent refractive indices that can simulate
  field-dependent $\mu$-crossovers.
\item Specific experimental configurations for measuring analog Hawking
  radiation in fiber systems are described, providing a direct
  experimental context for the DFD analog correspondence.
\end{enumerate}

\subsection{Updated Assessment}

The W4i13 GRIN fiber concept is now supported by established
experimental literature.  The specific enhancement needed for DFD
testing (a saturable nonlinearity implementing $\mu(x) = x/(1+x)$)
goes beyond standard Kerr nonlinearity, but semiconductor saturable
absorber mirrors (SESAMs) embedded in fiber could provide this.

\textbf{Cost estimate}: Unchanged at $<\$5$K for off-the-shelf
components.  The key experimental challenge is achieving sufficient
dynamic range to observe the crossover from the ``Newtonian'' (linear)
to ``MOND'' (saturated) regime within a single fiber length.


%=============================================================================
\part{Corpus Consistency Audit}
%=============================================================================

\section{Inconsistencies Found}

\begin{tcolorbox}[colback=red!5!white, colframe=red!60!black,
  title=\textbf{INCONSISTENCIES REQUIRING RESOLUTION}]

\begin{enumerate}

\item \textbf{Neutrino mass margin is STALE (Finding 1).}  The corpus
  states ``margin only 2.8~meV vs DESI DR2 bound.''  This refers to
  the Bayesian \LCDM{} bound of 72~meV from DESI DR1.  The current
  tightest bound is 53~meV (frequentist, DESI DR2 + Planck PR4 + CMB
  lensing).  Under this bound, DFD is \emph{excluded} at $>95\%$ CL.
  Under the Bayesian \LCDM{} bound of $\sim 72$~meV, DFD is still
  consistent.  The corpus should report \emph{both} bounds.

\item \textbf{$D_H/r_d$ sign-change table is UNRELIABLE (Finding 2).}
  The W4i12 table contains a known ``50$\times$ amplitude error'' and
  the sign at $z = 2.33$ contradicts the sign-change prediction.  This
  table should be marked as unreliable.  The W4i13 Survey Agent flagged
  this but no correction has been made.

\item \textbf{DESI CPL tension description is INTERNALLY INCONSISTENT.}
  The corpus states ``3.5--4.2$\sigma$'' tension (W4i13), but also
  says ``CPL is the WRONG observable for DFD'' (W4i11).  If CPL is
  the wrong observable, then the $\sigma$-level computed from CPL
  parameters should not be used to assess DFD tension.  The corpus
  should either (a) stop quoting CPL $\sigma$-levels as DFD tension,
  or (b) clearly label them as ``tension under wrong observable ---
  not a DFD constraint.''

\item \textbf{Falsification experiment count is inconsistent.}  The
  corpus lists ``5 falsification experiments'' (Section 4.7) and the
  context says ``MAQRO + PTA breathing mode = 6 falsification
  experiments.''  The W4i13 P1+P8 update recommended adding MAQRO as
  a 6th.  The corpus count should be updated to 6.

\end{enumerate}

\end{tcolorbox}


%=============================================================================
\part{Summary and Recommendations}
%=============================================================================

\begin{tcolorbox}[colback=green!5!white, colframe=green!60!black,
  title=\textbf{W4i14 Summary}]

\textbf{P1}: Confirmed DEAD (4th consecutive confirmation). No change.

\textbf{P8}: Confirmed DEAD (4th consecutive confirmation). No change.

\textbf{Five findings from audit and deep dive}:

\begin{enumerate}[nosep]
\item \textbf{Neutrino mass existential threat} (EXISTENTIAL):
  $\Sigma m_\nu = 61.4$~meV exceeds the tightest bound (53~meV at
  95\% CL, frequentist).  Model-dependent: relaxes to $\sim 120$~meV
  under $w_0 w_a$CDM.  DFD Boltzmann code CRITICAL to resolve.
  \textbf{Flag prediction as AT RISK.}

\item \textbf{$D_H/r_d$ table audit} (CRITICAL): W4i12 table has
  internal sign error at $z = 2.33$, contradicting the sign-change
  prediction.  Known 50$\times$ amplitude error unresolved.
  \textbf{Mark table as UNRELIABLE.}

\item \textbf{PPN preferred-frame parameter} (HIGH): $\alpha_1^{\rm DFD}
  \sim 10^{-13}$, eight orders below current bound $4\ee{-5}$.
  Undetectable but non-zero.  \textbf{New honest negative.}

\item \textbf{$S_8$ landscape shift} (HIGH): Some surveys now agree
  with Planck, weakening $S_8$ as a DFD discriminator.  DES Y6
  ``validation'' is conditional.  \textbf{Euclid October 2026 is
  decisive.}

\item \textbf{Fiber analog literature support} (MODERATE):
  arXiv:2512.15695 (Dec 2025) establishes the DFD-relevant optical
  analog framework.  \textbf{W4i13 GRIN concept now has direct
  literature backing.}
\end{enumerate}

\end{tcolorbox}

\section*{Action Items for Corpus Maintainer}

\begin{enumerate}
\item \textbf{URGENT}: Update neutrino mass entry from ``margin 2.8~meV
  vs DESI DR2 bound'' to ``excluded by 8.4~meV under frequentist
  \LCDM{} bound (53~meV); consistent under Bayesian \LCDM{} bound
  (72~meV); relaxed to $\sim 120$~meV under $w_0 w_a$CDM.  Status:
  AT RISK pending DFD Boltzmann code.''

\item \textbf{URGENT}: Mark W4i12 $D_H/r_d$ table as UNRELIABLE.
  Downgrade sign-change prediction to Tier 2 pending Boltzmann code.

\item \textbf{ADD}: CPL tension disclaimer --- either stop quoting
  CPL $\sigma$-levels as DFD tension, or clearly label as ``wrong
  observable.''

\item \textbf{ADD}: PPN $\alpha_1 \sim 10^{-13}$ as honest negative
  (preferred-frame effect, undetectable).

\item \textbf{ADD}: $S_8$ validation caveat --- conditional on tension
  being scale-dependent and persisting through Euclid.

\item \textbf{UPDATE}: Falsification experiment count from 5 to 6
  (MAQRO).

\item \textbf{ADD}: arXiv:2512.15695 as supporting reference for
  analog gravity / GRIN fiber concept.
\end{enumerate}


\end{document}
